\section{Dados e protocolo temporal da Bahia}
\label{sec:ba2026_methods}

A avaliação utiliza 145 séries loja-produto da Bahia, em 115 lojas e nove produtos, com 38 ciclos comerciais disponíveis. O objetivo é comparar uma política global com a escolha por perfil em períodos posteriores aos usados no ajuste e na seleção. A coorte foi definida retrospectivamente; os ajustes e a seleção, porém, utilizam apenas a informação disponível antes do teste. Essa distinção limita a generalização para séries ainda não observadas.

\subsection{Protocolo e cobertura}

Os registros complementares de vendas são somados por chave, sem descartar transações apenas por compartilharem loja, produto e ciclo. A coorte tem identificadores fixos; em cada origem, exige ao menos três ciclos positivos no treino. As 145 séries são elegíveis nos dois prefixos, de 17 e 23 ciclos. Cada prefixo é seguido por sete ciclos de validação e sete de teste. Na numeração iniciada em um, os testes abrangem os ciclos 25--31 e 31--37: compartilham um ciclo e representam medidas repetidas das mesmas séries.

Perfis, escalas dos estados e parâmetros são calculados no treino. O prazo efetivo de entrega é de dois ciclos, e EOQCycle utiliza demanda e manutenção na mesma unidade temporal. Custos de manutenção, falta, emissão e quantidade permanecem, respectivamente, em 1 por unidade e ciclo, 5 por unidade não atendida, 50 por pedido e 0,5 por unidade pedida. São unidades do cenário, sem equivalência monetária calibrada. O custo soma manutenção, falta e pedidos no bloco de teste, sem crédito pelo estoque final. Estados iniciais comuns são obtidos por uma política EOQ de referência: um aquecimento no passado precede a validação, e sua trajetória na validação determina o estado inicial do teste. O teste não continua o estoque particular de cada candidata.

O portfólio principal contém EOQCycle, sSCycle, CappedBaseStock, GA, PPO e GA--DQN; DQN é o controle adicional do híbrido. As três primeiras políticas têm um ajuste por série e origem. GA, PPO, GA--DQN e DQN usam sementes 11, 29 e 47, totalizando 4.350 ajustes e 8.700 registros de validação ou teste. A matriz está completa, sem casos excluídos por falta de política ou semente. O GA usa população de 100 indivíduos e 50 gerações. DQN e PPO têm orçamento mínimo de 20 mil transições, arredondado para episódios completos: 20.009 e 20.010 nas duas origens. GA--DQN acrescenta busca genética e mil transições de demonstração; seu esforço total é maior que o do DQN isolado. A variante avaliada mantém o mesmo gatilho e piso de pedido nas demonstrações, no treinamento e na inferência.

As sementes são primeiro promediadas dentro de cada série, origem e política. Para cada loja excluída, as políticas global e por perfil são escolhidas na validação das demais lojas e avaliadas no teste da loja reservada. O limiar principal continua em 0,70. Escolher uma alternativa viável na validação não garante repetir seu atendimento no futuro. Se nenhuma candidata cumpre a meta, a escolha de maior serviço é registrada como alternativa de contingência, sem classificá-la como viável.


\begin{figure}[htbp]
\centering
\includegraphics[width=\linewidth]{figures/fluxo_experimental.png}
\caption{Fluxo da avaliação e separação da informação. Cada candidata é ajustada no treino da própria série. A estratégia de seleção usa a validação das demais lojas e é avaliada no teste da loja reservada. Os dois testes compartilham um ciclo; as origens são medidas repetidas.}
\label{fig:experimental_flow}
\end{figure}

O fluxo da Figura~\ref{fig:experimental_flow} mantém separadas duas decisões: ajustar o controlador e escolher qual controlador recomendar. Para cada caso, o atendimento é $F=\sum_t v_t/\sum_t d_t$, em que $v_t$ é a demanda atendida e $d_t$ a demanda observada. Quando o denominador é zero, $F$ permanece indefinido. A média por caso atribui peso igual às séries com demanda; o atendimento ponderado usa a razão entre todas as unidades atendidas e demandadas. Ambos acompanham o custo total, para evitar que baixo atendimento seja interpretado como eficiência econômica.
